\documentclass{article}


\usepackage[utf8]{inputenc}
\usepackage[T1]{fontenc}
\usepackage[portuguese,brazil]{babel}
 
% Margens ampliadas
\usepackage[
  left=3cm,
  right=3cm,
  top=3cm,
  bottom=3cm
]{geometry}
 
% Matemática e símbolos
\usepackage{amsmath}
\usepackage{amssymb}
\usepackage{mathtools}
 
% Tabelas
\usepackage{booktabs}
\usepackage{array}
\usepackage{multirow}
\usepackage{colortbl}
 
% Gráficos e figuras
\usepackage{graphicx}
\usepackage{float}
 
% Cores
\usepackage{xcolor}
 
% Links e referências
\usepackage{hyperref}
\hypersetup{
  colorlinks=true,
  linkcolor=blue,
  urlcolor=blue,
  citecolor=blue,
  pdftitle={Proposta de Teste A/B - PhoneScorer},
  pdfauthor={},
  pdfsubject={Teste A/B}
}
 
% Espaçamento
\usepackage{setspace}


\title{Desafio Técnico - Cientista de Dados Pleno: \\Proposta Teste A/B}
\author{Maria Luiza Wuillaume}
\date{26/04/2026}

\begin{document}

\maketitle

\section{Contexto e Motivação}

O algoritmo \texttt{PhoneScorer} propõe um ranking interpretável baseado em três sinais
(confiabilidade da fonte, atualidade e histórico de engajamento) em substituição às estratégias
atualmente em uso (envio aleatório ou ordem alfabética).

Para avaliar seu impacto, definimos um experimento A/B no qual:
\begin{itemize}
    \item o \textbf{grupo controle} utiliza a estratégia atual de seleção de telefones (aleatório ou ordem alfabética);
    \item o \textbf{grupo tratamento} utiliza o ranking gerado pelo \texttt{PhoneScorer}.
\end{itemize}

O objetivo é verificar se a nova abordagem produz ganhos reais no motor de disparos de WhatsApp,
com base em métricas de engajamento bem definidas ao longo de um período controlado.
 
 \section{Métricas}

\subsection{Métrica Primária}

\textbf{Proporção de CPFs com alta taxa de leitura:}

A taxa de leitura por CPF é definida como:

\[
\text{ReadRate}_{\text{CPF}} = \frac{\text{Total de mensagens lidas}}{\text{Total de mensagens enviadas}}
\]

Definimos a variável indicadora:

\[
\text{HighRead}_{\text{CPF}} =
\begin{cases}
1, & \text{se } \text{ReadRate}_{\text{CPF}} \geq 0.75 \\
0, & \text{caso contrário}
\end{cases}
\]

A métrica primária é:

\[
p_{\text{read}} = \frac{\sum \text{HighRead}_{\text{CPF}}}{n}
\]

onde $n$ é o número total de CPFs considerados no experimento.

\vspace{0.3cm}

\subsection{Métrica Secundária}

\textbf{Proporção de CPFs com alta taxa de entrega:}

A taxa de entrega por CPF é definida como:

\[
\text{DeliveryRate}_{\text{CPF}} = \frac{\text{Mensagens com status delivered ou read}}{\text{Total de mensagens enviadas}}
\]

Definimos a variável indicadora:

\[
\text{HighDelivery}_{\text{CPF}} =
\begin{cases}
1, & \text{se } \text{DeliveryRate}_{\text{CPF}} \geq 0.9 \\
0, & \text{caso contrário}
\end{cases}
\]

A métrica secundária é:

\[
p_{\text{delivery}} = \frac{\sum \text{HighDelivery}_{\text{CPF}}}{n}
\]



 \section{Hipóteses}

\subsection{Hipótese Nula ($H_0$)}

\[
H_0: p_{\text{read, tratamento}} \leq p_{\text{read, controle}}
\]

A hipótese nula afirma que o grupo tratamento (\texttt{PhoneScorer}) não aumenta a proporção de CPFs com alta taxa de leitura em relação ao grupo controle.

\vspace{0.3cm}

\subsection{Hipótese Alternativa ($H_1$)}

\[
H_1: p_{\text{read, tratamento}} > p_{\text{read, controle}}
\]

A hipótese alternativa afirma que o grupo tratamento (\texttt{PhoneScorer}) aumenta a proporção de CPFs com alta taxa de leitura em relação ao grupo controle.

\section{Teste Estatístico}

A comparação entre os grupos é realizada por meio de um \textbf{teste de duas proporções (z-test)}, adequado para avaliar diferenças entre proporções em amostras independentes.

O teste é conduzido de forma unilateral à direita, com nível de significância $\alpha = 0.05$, avaliando se:

\[
p_{\text{read, tratamento}} > p_{\text{read, controle}}
\]

Rejeita-se $H_0$ se o $p$-value $< 0.05$.

A decisão final considera evidência de melhoria na métrica primária sem degradação relevante na métrica secundária.




\section{Cálculo de Tamanho de Amostra}

O tamanho de amostra necessário para o experimento foi estimado com base na métrica primária (proporção de CPFs com alta taxa de leitura).

Para o cálculo, foi utilizada a calculadora de Evan Miller:

\url{https://www.evanmiller.org/ab-testing/sample-size.html}

Os parâmetros foram definidos a partir da análise exploratória realizada no notebook \url{07_proporcao_cpfs_taxas.ipynb}.

\begin{itemize}
    \item \textbf{Baseline (taxa atual)}: $70\%$, correspondente à proporção observada de CPFs com taxa de leitura $\geq 90\%$ (69.76\% em 268{,}436 CPFs analisados).

    \item \textbf{Minimum Detectable Effect (MDE)}: $5$ pontos percentuais, representando um aumento de $70\%$ para $75\%$.

    \item \textbf{Nível de significância}: $\alpha = 0.05$

    \item \textbf{Poder estatístico}: $1 - \beta = 0.80$
\end{itemize}

Com esses parâmetros, o tamanho de amostra estimado foi de aproximadamente:

\[
n \approx 1{,}335 \text{ CPFs por grupo}
\]




\section{Duração do Experimento}

A duração do experimento foi definida considerando o tamanho de amostra necessário e a variação temporal no comportamento dos usuários. Recomenda-se uma duração mínima de \textbf{7 dias consecutivos}, suficiente para capturar um ciclo semanal completo (dias úteis e finais de semana).

Com base no tamanho de amostra estimado de aproximadamente $2{,}670$ CPFs no total, é necessário observar em média cerca de:

\[
\approx 380 \text{ CPFs por dia}
\]

para atingir o volume necessário ao longo de 7 dias.

Adicionalmente, para cada mensagem enviada, considera-se uma janela de \textbf{24 horas} para observação de leitura. Dessa forma, as métricas de leitura e entrega são sempre calculadas considerando esse mesmo horizonte temporal, assegurando consistência entre os grupos.


\section{Design do Experimento}

\begin{table}[h]
\centering
\small
\begin{tabular}{|l|c|c|}
\hline
\textbf{Aspecto} & \textbf{Grupo Controle} & \textbf{Grupo Tratamento} \\
\hline
\textit{Estratégia de Seleção} & Aleatória ou alfabética (atual) & PhoneScorer \\
\hline
\textit{Critério de Alocação} & \multicolumn{2}{c|}{Aleatório por CPF (50/50)} \\
\hline
\textit{Tamanho} & \multicolumn{2}{c|}{1.335 CPFs por grupo} \\
\hline
\textit{Período} & \multicolumn{2}{c|}{7 dias consecutivos} \\
\hline
\end{tabular}
\caption{Estrutura do Teste A/B: alocação aleatória de CPFs a controle e tratamento.}
\label{tab:design_ab}
\end{table}

\subsection*{Definições do Experimento}

\begin{itemize}
    \item \textbf{Unidade de Randomização:} CPF (não telefone). Cada CPF é alocado uma única vez a um dos grupos no início do experimento.

    \item \textbf{Persistência da Alocação:} A alocação é fixa ao longo de todo o experimento (\textit{sticky assignment}), garantindo que um mesmo CPF não transite entre controle e tratamento.

    \item \textbf{Unidade de Análise:} As métricas são agregadas no nível de CPF.

    \item \textbf{Regra de Exposição:} Para cada CPF, o grupo tratamento (PhoneScorer) seleciona o \textbf{telefone com maior score} (top-1) para envio da mensagem. O grupo controle segue a estratégia atual (aleatória ou alfabética).

    \item \textbf{Janela de Observação:} Para cada mensagem enviada, o status de entrega e leitura é avaliado após uma janela fixa de 24 horas.

    \item \textbf{Critérios de Inclusão:} A análise considera apenas mensagens com janela de observação completa de 24 horas. CPFs sem envios no período não são incluídos no cálculo das métricas.

    \item \textbf{Balanceamento:} A alocação é realizada de forma balanceada entre os grupos (50/50), resultando em aproximadamente 1.335 CPFs por grupo.

    \item \textbf{Independência:} Cada CPF é tratado como uma observação independente no experimento.
\end{itemize}
\end{document}
